% Source: https://artofproblemsolving.com/community/c3h2373987_all_the_pathways_problem
% Stable identity: aops-topic-2373987-q01
A spelling is a sequence of squares $P_1,I,P_3,U,P_5$, each next to the previous one. All P's are on the outer ring; the I's are the four inner corner squares and the U's are the four inner edge squares.

\textbf{Neighbours.} Number the rows and columns $1$ to $5$ from the top left, and take the I at $(2,2)$. Its P-neighbours are the five outer squares $(1,1),(1,2),(1,3),(2,1),(3,1)$. The U at $(2,3)$ has three P-neighbours, $(1,2),(1,3),(1,4)$, and two of them, $(1,2)$ and $(1,3)$, are also next to this I. Likewise the U at $(3,2)$ has three P-neighbours, $(2,1),(3,1),(4,1)$, two of which are next to this I. The other two U's, at $(3,4)$ and $(4,3)$, have their P-neighbours on the right and bottom edges only, none of them next to this I. By symmetry the same holds for every I.

\textbf{Counting.} Fix the I and the U. Then $P_1$ is any P next to the I ($5$ ways), $P_3$ is any P next to both the I and the U ($2$ ways), and $P_5$ is any P next to the U ($3$ ways), and these choices are independent because squares may be reused. So each suitable (I, U) pair gives $5\cdot2\cdot3=30$ spellings. Each of the $4$ I's has exactly $2$ suitable U's (a U sharing no P-neighbour with the I gives no spelling), and the total is
\[4\cdot2\cdot30=\boxed{240}.\]
